\documentclass[11pt]{article}

% ---------------------------------------------------------------
%  AI Applications in Research: Concepts and the Use of Claude Code
%  A conceptual summary, informed by Scott Cunningham's
%  "Scott's Mixtape Substack" Claude Code series (2026).
% ---------------------------------------------------------------

\usepackage[utf8]{inputenc}
\usepackage[T1]{fontenc}
\usepackage[margin=1in]{geometry}
\usepackage{amsmath}
\usepackage{enumitem}
\usepackage{hyperref}
\usepackage{xcolor}
\usepackage{titlesec}

\hypersetup{
  colorlinks=true,
  linkcolor=blue!55!black,
  urlcolor=blue!55!black,
  citecolor=blue!55!black
}

\titleformat{\section}{\large\bfseries}{\thesection}{0.6em}{}
\titleformat{\subsection}{\normalsize\bfseries}{\thesubsection}{0.6em}{}

\newcommand{\concept}[1]{\textbf{#1}}

\title{\textbf{AI Applications in Research}\\[0.3em]
\large Core Concepts and the Use of Claude Code}
\author{Prepared for M.\ Caruso}
\date{14 June 2026}

\begin{document}

\maketitle

\begin{abstract}
\noindent This document summarizes the main concepts surrounding the use of
artificial intelligence (AI) in scientific research, with particular emphasis
on agentic coding assistants and \concept{Claude Code}. It synthesizes general
principles of AI-assisted research with field observations drawn from Scott
Cunningham's ongoing ``Claude Code'' series on \emph{Scott's Mixtape Substack},
in which an applied econometrician documents the practical promise and the new
failure modes of conducting empirical research with large language model (LLM)
agents. The central thesis is that AI agents \emph{solve old problems and create
new ones}, and that the researcher's job is to build workflows that solve the
new ones---above all, to guarantee verification and ``zero error.''
\end{abstract}

\section{Introduction}
AI---and in particular LLM-based \concept{agents}---has moved from a writing aid
to an active participant in the research pipeline. Where earlier tools
autocompleted code or text, agentic systems such as Claude Code can read a
project, plan multi-step tasks, write and execute scripts, produce figures and
tables, run version control, and reason about results in natural-language
dialogue. This shift changes not only \emph{how fast} research is produced but
also \emph{where} errors originate and how they are caught. This document lays
out the main concepts a researcher should understand to use such tools well.

\section{What Is Claude Code?}
\concept{Claude Code} is Anthropic's command-line agent built on the Claude
family of models. Rather than returning isolated snippets, it operates inside a
project directory and can:
\begin{itemize}[leftmargin=1.4em,itemsep=2pt]
  \item read and edit files across a codebase while preserving structure;
  \item run shell commands, execute analysis scripts, and inspect their output;
  \item produce replicable code, organized directories, consistent naming
        conventions, and automated tables/figures;
  \item perform \concept{git} operations, making version control routine;
  \item hold an extended, conversational \concept{dialogue} about the task,
        acting as a co-author of the research process rather than a one-shot
        code generator.
\end{itemize}
The most capable current Claude models include the Opus, Sonnet, and Haiku
tiers; agentic research workflows typically favor the most capable tier for
reasoning-heavy steps.

\section{Main Concepts in AI-Assisted Research}

\subsection{The research pipeline as the unit of automation}
The natural object to automate is not a single script but the whole
\concept{pipeline}: data ingestion, cleaning, analysis, exhibit production, and
write-up. Claude Code is attractive precisely because it internalizes good
software-engineering hygiene (modular scripts, reproducible builds, version
control) that researchers historically had to impose on themselves by hand.

\subsection{Verification is separate from production}
A defining property of agents is that they \concept{cleanly separate production
from verification}: the agent produces the result, and the human is left to
verify it. In traditional hand-coding, production and verification were fused---
the act of writing the code generated an ``epistemological feeling'' of knowing
\emph{what made what}. With agents, that felt confidence is lost and must be
deliberately reconstructed.

\subsection{The returns to domain knowledge \emph{rise}}
A counterintuitive but central lesson: automation does \emph{not} reduce the
value of expertise---it raises it. Because the agent can silently make
\concept{reasoning mistakes} (not just coding bugs), only a researcher with deep
domain ``human capital'' can recognize when an output is implausible. Cunningham
frames this as: \emph{the returns to econometrics knowledge with Claude Code are
higher, not lower.}

\subsection{``Zero error'' as the binding constraint}
The recommended mental model is to treat \concept{zero error} as a
non-negotiable constraint and to make everything else---which skills, which
prompts, which scaffolding---an endogenous choice variable optimized to satisfy
it. Crucially, when errors slip through, it is the \emph{researcher}, not the
agent, who bears the professional cost.

\subsection{Estimand-first (``forward engineering'')}
Before letting an agent compute anything, the researcher should write down the
target \concept{estimand} explicitly---in potential-outcomes notation, naming
the population and the weighting of the averaged effects. This ``forward
engineering'' discipline (e.g., the practitioner's-guide approach to
difference-in-differences) prevents the agent from quietly
\concept{specification searching} and reporting only the method that happens to
fit the prompt.

\subsection{Structured (and analog) workflows over ad hoc skills}
Rather than relying on packaged \concept{/skills} and ``plan-then-let-it-rip''
automation, the field experience favors \concept{structured workflows}---often
explicit \concept{checklists} executed step by step, where one does not proceed
until the prior step is verified. The analogy is Atul Gawande's
\emph{The Checklist Manifesto}: checklists cut errors in surgery and aviation
not by adding intelligence but by enforcing discipline.

\section{New Failure Modes Introduced by Agents}
AI agents create genuinely new categories of error that look different from
human mistakes and are therefore harder to detect:
\begin{description}[leftmargin=1.4em,style=nextline,itemsep=3pt]
  \item[Drift] The analytical sample, treatment definitions, or scripts shift
    silently across sessions; the foundation moves without an obvious signal.
  \item[Unwritten work] Steps performed ``in dialogue'' may never be committed
    to a script, so the pipeline cannot be reproduced.
  \item[A ``Claude mess''] Output that is sprawling yet superficially tidy---a
    well-designed sprawl that does not \emph{look} like a mess, which makes
    stale or duplicated artifacts hard to spot.
  \item[Hidden decisions] Agents may store many under-the-hood choices (e.g.,
    in a JSON state file) that the researcher never sees unless they look.
  \item[Confident wrong reasoning] Agents will happily pursue and ``conclude''
    on faulty conjectures (e.g., debugging for hours while the real cause is a
    wrong package version), with no internal sense that the conclusion is wrong.
  \item[Amnesia] Because the human did not produce the scaffolding, returning to
    a project after even a couple of days leaves the researcher unusually
    disoriented, lacking the memorized mental map that hand-production created.
\end{description}

\section{Practical Strategies and Mitigations}
The following practices recur as antidotes to the failure modes above:
\begin{enumerate}[leftmargin=1.5em,itemsep=3pt]
  \item \concept{Stay in the loop.} Do not step away and trust the plan; verify
    continuously, because automation is not verification.
  \item \concept{Pre-register the estimand} in potential-outcomes terms before
    any estimation.
  \item \concept{Use a domain checklist} (e.g., name the estimand; tabulate
    treated units by cohort; visualize rollout; justify covariates; plot outcome
    means; run honest intervals and falsification tests) and complete each step
    before advancing.
  \item \concept{Build a harness / dashboard} that captures the workflow, makes
    figures and tables clickable ``cards'' revealing the code that produced
    them, and supports \concept{psychological verification}---rebuilding the felt
    confidence that hand-coding used to supply.
  \item \concept{Maintain progress logs} after each session and auto-generate a
    concise ``deck'' that re-orients you on return, directly countering amnesia.
  \item \concept{Keep an analog map.} A physical notebook to track what the agent
    is doing constrains its latitude and anchors human memory.
  \item \concept{Design systems endogenous to your own constraints} (working
    style, strengths, even conditions such as ADHD) rather than assuming an
    idealized operator; the right verification method is personal and not
    transferable wholesale from another researcher.
\end{enumerate}

\section{A Worked Domain: Staggered Difference-in-Differences}
Cunningham's running example is the Callaway and Sant'Anna (CSDID) estimator,
where the estimator reduces conceptually to ``four averages and three
subtractions'' with weights on cell means. The lesson generalizes: an expert can
verify an agent's output \emph{only} because they know the method's structure---
e.g., that there should be no coefficient at the period $g-1$ baseline, that
$7$--$10$ ``events'' per covariate are needed to fit a propensity score, or that
a normalized mean difference
\[
  \mathrm{ND} = \frac{\bar X_1 - \bar X_0}{\sqrt{\tfrac{1}{2}\left(\mathrm{var}_1 + \mathrm{var}_0\right)}}
\]
above $0.25$ in absolute value flags a covariate that must be included to defend
conditional parallel trends. Without such human capital, an agent's plausible-
looking---but wrong---output is indistinguishable from a correct one.

\section{Conclusion}
AI agents like Claude Code are best understood not as a replacement for the
researcher but as a powerful collaborator that \emph{relocates} where effort and
risk live. They absorb much of the mechanical burden of coding and reproducibility
while shifting the researcher's role toward \concept{specification},
\concept{verification}, and \concept{workflow design}. The recurring, hard-won
conclusions are: domain expertise becomes \emph{more} valuable, verification must
be rebuilt deliberately because the agent severs it from production, and ``zero
error'' should be the constraint around which every other choice is optimized.
Used with that discipline, AI is a genuine accelerant for research; used without
it, it quietly manufactures confident, hard-to-detect mistakes.

\vfill
\noindent\rule{\linewidth}{0.4pt}\\[0.3em]
\footnotesize Sources and inspiration: Scott Cunningham, ``Claude Code'' series
(esp.\ \#53 ``Applied econometrics will require a detailed checklist'' and \#54
``Amnesia'') and related posts, \emph{Scott's Mixtape Substack},
\href{https://causalinf.substack.com}{causalinf.substack.com}, May--June 2026;
Atul Gawande, \emph{The Checklist Manifesto}; Baker et al., ``Difference-in-
Differences: A Practitioner's Guide,'' \emph{Journal of Economic Literature},
2026.

\end{document}
